\begin{cabstract}
应概括地反映出本论文的主要内容，具有独立性和自创性，包括研究工作目的、方法、结果和结论，要突出本论文的创造性成果。摘要力求语言精炼准确，硕士学位论文建议1000字以内，博士学位论文建议2000字以内。部分学生用外文撰写学位论文时，博士学位论文的摘要应不少于5000字符，硕士应不少于2500字符。摘要中不可出现图片、图表、表格或其他插图材料。

关键词是为了便于做文献索引和检索工作而从论文中选取出来用以表示全文主题内容信息的单词或术语。关键词应体现论文特色，具有语义性，在论文中有明确的出处，并应尽量采用《汉语主题词表》或各专业主题词表提供的规范词。在摘要内容下方另起一行标明，一般3～8个，之间用‘；’分开。为便于国际交流，应标注与中文对应的英文关键词。
\end{cabstract}
\ckeywords{关键词1,关键词2,关键词3,关键词4,关键词5}
\begin{eabstract} 
The abstract content corresponds to the Chinese abstract. It should be noted that the case of foreign language and the writing in italics should meet the relevant requirements.

Abstract内容与中文摘要相对应，应注意外文大小写、正斜体书写符合有关要求。
\end{eabstract}

\ekeywords{Keywords1,Keywords2,Keywords3,Keywords4,Keywords5}
